\subsection{Discusión}
\label{sec:discusion-pruebas}

\subsubsection{Hallazgos principales}
\begin{itemize}
    \item \textbf{Aceptación por parte de los alumnos.} El agrado general y la disposición a volver a practicar fueron altos (16 de 17 alumnos en ambas preguntas). Los apoyos multimodales de la vista del alumno, es decir, los dibujos y la lectura en voz alta de las instrucciones, fueron valorados positivamente.
    \item \textbf{La escritura en pantalla es el principal punto de fricción.} Fue la pregunta con menos respuestas afirmativas (10 de 17) y la que reunió más respuestas indecisas. Este resultado concuerda con la valoración docente del tamaño de textos e íconos y de la autonomía del alumno. En conjunto, apunta a revisar el tamaño del lienzo, el grosor del trazo y el tamaño de los controles de la vista del alumno.
    \item \textbf{La revisión docente fue valorada.} Ambos docentes consideraron fácil de entender la calificación sugerida y adecuado que el maestro la confirme o corrija antes de que cuente. Esto respalda la decisión de diseño de que las sugerencias de CLS y legibilidad solo se incorporen al progreso tras la confirmación del tutor.
    \item \textbf{Complemento del cuaderno.} La mayoría de los alumnos prefirió escribir tanto en la tableta como en el cuaderno, lo que coincide con el enfoque de la aplicación como herramienta complementaria.
    \item \textbf{Barreras de infraestructura.} Los dos docentes identificaron la falta de internet y de dispositivos como dificultades para el uso en su escuela. Son condiciones externas al sistema, pero limitan su adopción.
\end{itemize}

\subsubsection{Mejoras derivadas}
A partir de estos resultados se identifican las siguientes mejoras:
\begin{enumerate}
    \item Aumentar el tamaño de la letra, de los íconos y de los botones de la vista del alumno, y revisar la superficie y el grosor del trazo del lienzo.
    \item Reforzar las instrucciones iniciales de cada ejercicio, por ejemplo con un ejemplo animado o la reproducción automática del audio, para reducir las respuestas \textit{Más o menos} sobre la comprensión de la actividad.
    \item Elaborar una guía breve o una sesión de capacitación para docentes, como lo sugirió uno de ellos.
\end{enumerate}

\subsubsection{Limitaciones}
\begin{itemize}
    \item La muestra es pequeña y por conveniencia: 17 alumnos y solo 2 docentes. Los resultados no pueden generalizarse y, en el caso de los docentes, solo describen la opinión de los participantes.
    \item Se trata de una sesión única, sin grupo de control ni mediciones antes y después, por lo que no permite atribuir cambios en el desempeño a la aplicación. El seguimiento de un mes previsto para la prueba piloto sigue pendiente.
    \item Los niños pequeños tienden a responder de forma positiva en las escalas de satisfacción \cite{read2006funtoolkit}. Además, las preguntas fueron leídas en voz alta por el docente, lo que pudo influir en las respuestas.
    \item La pregunta abierta de los alumnos obtuvo una respuesta casi idéntica en la mayoría de los casos, por lo que aporta poca información sobre dificultades específicas.
    \item Las encuestas miden percepción y aceptación. No incluyen todavía la comparación cuantitativa entre las sugerencias del sistema (CLS y medidas de legibilidad) y la valoración docente, ni la tasa de error del reconocimiento (CER y WER), descritas en la Sección~\ref{sec:validacion-vision}. % TODO-CONFIRMAR: si existen estos datos de la sesión, incorporarlos aquí.
\end{itemize}
